\chapter{Droughts}
\index{Droughts}

\section*{Definition and Overview}
Drought is a prolonged period of unusually low rainfall that leads to water shortages, affecting communities, agriculture, and the environment. In many countries, droughts are recurring natural events with significant health and wellbeing effects, especially in rural and remote communities. Drought affects mental health through financial stress, isolation, and uncertainty, and it can affect physical health through hardship, poor nutrition, and increased workloads. While droughts cannot be prevented, their health effects can be reduced with support, planning, and community resilience. This chapter explains the health impacts of drought and how to manage them.

\section*{Epidemiology}
the affected region experiences droughts regularly, with major events such as the Federation, World War II, and Millennium droughts causing widespread hardship. Rural and farming communities bear the greatest burden, and farmers have significantly higher rates of psychological distress and suicide than the general population. Drought affects people indirectly too, through higher food prices, dust storms, bushfire risk, and water restrictions. Climate change is projected to increase the frequency and severity of drought in parts of some countries.

\section*{Aetiology and Pathophysiology}
\begin{itemize}
    \item \textbf{Water shortage:} prolonged low rainfall reduces water for drinking, farming, and firefighting
    \item \textbf{Financial stress:} crop and livestock losses reduce income and increase debt
    \item \textbf{Mental health effects:} stress, anxiety, depression, and increased suicide risk in affected communities
    \item \textbf{Physical effects:} heat, dust, and increased physical work affect health
    \item \textbf{Social effects:} isolation, population decline, and pressure on services
    \item \textbf{Environmental effects:} dust storms, air pollution, and reduced water quality
    \item \textbf{Mechanism:} prolonged stress and hardship affect both psychological and physical wellbeing
\end{itemize}

\section*{Clinical Features}
\begin{itemize}
    \item Stress, anxiety, and low mood related to financial and environmental pressures
    \item Sleep disturbance, irritability, and difficulty concentrating
    \item Physical symptoms such as fatigue, headaches, and respiratory problems from dust
    \item Increased alcohol use and family conflict
    \item Reduced access to services in affected regions
    \item Heat-related illness during drought heatwaves
    \item Suicide risk in the most affected groups
\end{itemize}

\section*{Differential Diagnosis}
\begin{itemize}
    \item Depression and anxiety from drought stress versus other causes
    \item The health effects of drought versus those of other disasters such as floods and fires
    \item Physical symptoms from dust and heat versus from other illnesses
\end{itemize}

\section*{When to Seek Medical Care}
People in drought-affected areas who experience persistent distress, sleep problems, or thoughts of harming themselves should seek help from a GP, mental health service, or support line. Financial and practical support services are available through rural support organisations.

\textbf{Contact your local emergency services for:}
\begin{itemize}
    \item Thoughts of self-harm or suicide
    \item Severe distress or crisis
    \item Heat stroke or severe dehydration
    \item Any medical emergency
\end{itemize}
\textbf{Support lines:} a local crisis service and the Rural and Remote Mental Health support services.

\section*{Diagnosis and Investigations}
\begin{itemize}
    \item \textbf{Psychological assessment:} mood, stress, and coping
    \item \textbf{Risk assessment:} suicide risk, substance use, and family violence
    \item \textbf{Physical assessment:} dust-related respiratory symptoms, heat illness, and general health
    \item \textbf{Practical assessment:} financial and service access needs
    \item \textbf{Referral:} to rural health services, counselling, and support organisations
\end{itemize}

\section*{Management}
\begin{itemize}
    \item \textbf{Mental health support:} counselling, telehealth services, and peer support
    \item \textbf{Financial support:} farm support programs, drought assistance, and financial counselling
    \item \textbf{Practical support:} rural support organisations and community services
    \item \textbf{Health care:} treatment of anxiety, depression, and physical conditions
    \item \textbf{Community resilience:} local networks, social connection, and planning
    \item \textbf{Water management:} household and farm water planning
    \item \textbf{Dust and heat precautions:} protecting health during dust storms and heatwaves
\end{itemize}

\section*{Complications}
\begin{itemize}
    \item Depression, anxiety, and suicide
    \item Financial ruin and loss of farms and livelihoods
    \item Family breakdown and violence
    \item Respiratory problems from dust
    \item Heat-related illness
    \item Depopulation and loss of services in affected communities
\end{itemize}

\section*{Prognosis and Outcomes}
Droughts end, and communities that maintain social connection, seek support, and plan ahead recover better. Mental health treatment is effective, and most people affected by drought stress recover with appropriate support. Government and community programs provide practical assistance during and after drought. Building resilience before and during drought reduces its health effects and supports recovery.

\section*{Prevention and Self-Care}
\begin{itemize}
    \item Maintain social connections and check on neighbours
    \item Seek financial and practical support early
    \item Talk about stress and seek counselling when needed
    \item Limit alcohol during stressful periods
    \item Protect against dust with masks and indoor air quality measures
    \item Prepare for heatwaves with hydration and cooling plans
    \item Use telehealth services when distance is a barrier
    \item Plan for the long term with diversified income where possible
\end{itemize}

\section*{Sources and Further Reading}
The following reputable sources provide further information for healthcare students and patients seeking reliable, current advice about drought and health.

\begin{itemize}
    \item National Centre for Farmer Health. \textit{Drought and farmer health.} https://www.farmerhealth.org.au/
    \item Rural Aid. \textit{Support for farming communities.} https://www.ruralaid.org.au/
    \item Beyond Blue. \textit{Mental health support for rural Australians.} https://www.beyondblue.org.au/
    \item Lifeline Australia. \textit{Support and crisis services.} https://www.lifeline.org.au/
    \item Australian Government Department of Agriculture. \textit{Drought assistance programs.} https://www.agriculture.gov.au/
    \item Healthdirect Australia. \textit{Managing health during drought.} https://www.healthdirect.gov.au/
\end{itemize}
